% questions.tex —— 2009 年·普通高等学校招生考试·新课标卷（琼、宁卷）·文科 · 题目内容片段（无答案）
% 由 tools/generate_projects.py 依来源总账自 2009/2009_数学真题_新课标文_试题.md 生成，勿手改（改源 MD 后重跑）


状态：needs-review；依固定候选 PDF 两页（印刷页 543–544）逐题核验，图形已矢量重绘，正式构建与逐页验收见回执。源未标分值，不补分值。原页第 22–24 题紧接第 21 题排印，未另印「四、选做题」标题；本库按 2008 同版式补出节标题并在此注明。原页若干排印问题照录并登记问题条目：第 9 题「棱线长」、第 19 题「抽插结果」、第 22 题第 (1) 问重复印出「证明：」均照录；第 16 题图象只画出 $x\geqslant\frac{\pi}{4}$ 的部分，照原图重绘并登记。

\gkpart{一、选择题}

\begin{question}
\gkti{1}已知集合 $A=\{1,3,5,7,9\}$，$B=\{0,3,6,9,12\}$，则 $A\cap B=$（　）

（A）$\{3,5\}$

（B）$\{3,6\}$

（C）$\{3,7\}$

（D）$\{3,9\}$

\end{question}
\begin{question}
\gkti{2}复数 $\frac{3+2\mathrm i}{2-3\mathrm i}=$（　）

（A）$1$

（B）$-1$

（C）$\mathrm i$

（D）$-\mathrm i$

\Needspace{0.4\textheight}
\end{question}
\begin{question}
\gkti{3}对变量 $x$，$y$ 有观测数据 $(x_i,y_i)$（$i=1,2,\cdots,10$），得散点图 1；对变量 $u$，$v$ 有观测数据 $(u_i,v_i)$（$i=1,2,\cdots,10$），得散点图 2。由这两个散点图可以判断（　）
\nopagebreak[4]
\begin{center}% 2009 年新课标卷文科第 3 题：两个变量的散点图（图 1：x-y；图 2：u-v）。
% 来源 PDF：普通高考/2009/2009新课标文(琼,宁).pdf，
% SHA256 9727691ae735985e787f7ae3febf48dc9e6c6cdd1aa53b465f5d95648abdd627
% 原页：第 1 页（A3 横向）左栏中部，印刷页 543。
% 重绘依据：图 1 的 10 个点整体随 x 增大而下降；图 2 的 10 个点整体随 u 增大而上升。
% 坐标按原图刻度估读，仅保留趋势，不改变题目所问的“正相关/负相关”判断。
% 纵向坐标按原数值除以 10 作图，刻度标注仍写原数值。
\begin{tikzpicture}[scale=0.82, semithick, >=Stealth,
  every node/.style={inner sep=1.5pt, font=\small}]
  % 图 1
  \begin{scope}
    \draw[->] (0,0) -- (7.6,0) node[right] {$x$};
    \draw[->] (0,0) -- (0,3.35) node[above] {$y$};
    \foreach \x in {1,...,7} \draw (\x,0) -- (\x,-0.12) node[below] {$\x$};
    \foreach \y in {5,10,15,20,25,30} {
      \draw (0,\y/10) -- (-0.14,\y/10) node[left] {$\y$};
    }
    \node[below left] at (0,0) {$O$};
    \foreach \x/\y in {0.8/23,1.4/20,2.0/20.5,2.4/16.5,3.2/15,
                      3.7/10.8,4.3/11.2,4.9/8,5.5/5.5,6.4/5}
      \fill (\x,\y/10) circle (2.2pt);
    \node at (3.8,-0.95) {图 1};
  \end{scope}
  % 图 2
  \begin{scope}[shift={(10.6,0)}]
    \draw[->] (0,0) -- (7.6,0) node[right] {$u$};
    \draw[->] (0,0) -- (0,6.6) node[above] {$v$};
    \foreach \x in {1,...,7} \draw (\x,0) -- (\x,-0.12) node[below] {$\x$};
    \foreach \y in {10,20,30,40,50,60} {
      \draw (0,\y/10) -- (-0.14,\y/10) node[left] {$\y$};
    }
    \node[below left] at (0,0) {$O$};
    \foreach \x/\y in {1.2/12,2.0/20,2.6/22,3.2/30,3.8/38,
                      4.2/36,4.6/46,5.2/41,5.9/54,6.5/57}
      \fill (\x,\y/10) circle (2.2pt);
    \node at (3.8,-0.95) {图 2};
  \end{scope}
\end{tikzpicture}
\end{center}
\nopagebreak[4]
（A）变量 $x$ 与 $y$ 正相关，$u$ 与 $v$ 正相关

（B）变量 $x$ 与 $y$ 正相关，$u$ 与 $v$ 负相关

（C）变量 $x$ 与 $y$ 负相关，$u$ 与 $v$ 正相关

（D）变量 $x$ 与 $y$ 负相关，$u$ 与 $v$ 负相关

\end{question}
\begin{question}
\gkti{4}有四个关于三角函数的命题：

$p_1$：$\exists x\in\mathbf R$，$\sin^2\frac x2+\cos^2\frac x2=\frac12$；

$p_2$：$\exists x,y\in\mathbf R$，$\sin(x-y)=\sin x-\sin y$；

$p_3$：$\forall x\in[0,\pi]$，$\sqrt{\frac{1-\cos2x}{2}}=\sin x$；

$p_4$：$\sin x=\cos y\Rightarrow x+y=\frac{\pi}{2}$。

其中假命题的是（　）

（A）$p_1$，$p_4$

（B）$p_2$，$p_4$

（C）$p_1$，$p_3$

（D）$p_2$，$p_3$

\end{question}
\begin{question}
\gkti{5}已知圆 $C_1:(x+1)^2+(y-1)^2=1$，圆 $C_2$ 与圆 $C_1$ 关于直线 $x-y-1=0$ 对称，则圆 $C_2$ 的方程为（　）

（A）$(x+2)^2+(y-2)^2=1$

（B）$(x-2)^2+(y+2)^2=1$

（C）$(x+2)^2+(y+2)^2=1$

（D）$(x-2)^2+(y-2)^2=1$

\end{question}
\begin{question}
\gkti{6}设 $x$，$y$ 满足 $\begin{cases}2x+y\geqslant4,\\ x-y\geqslant-1,\\ x-2y\leqslant2,\end{cases}$ 则 $z=x+y$（　）

（A）有最小值 $2$，最大值 $3$

（B）有最小值 $2$，无最大值

（C）有最大值 $3$，无最小值

（D）既无最小值，也无最大值

\end{question}
\begin{question}
\gkti{7}已知 $\boldsymbol a=(-3,2)$，$\boldsymbol b=(-1,0)$，向量 $\lambda\boldsymbol a+\boldsymbol b$ 与 $\boldsymbol a-2\boldsymbol b$ 垂直，则实数 $\lambda$ 的值为（　）

（A）$-\frac17$

（B）$\frac17$

（C）$-\frac16$

（D）$\frac16$

\end{question}
\begin{question}
\gkti{8}等差数列 $\{a_n\}$ 前 $n$ 项和为 $S_n$。已知 $a_{m-1}+a_{m+1}-a_m^2=0$，$S_{2m-1}=38$，则 $m=$（　）

（A）$38$

（B）$20$

（C）$10$

（D）$9$

\Needspace{0.4\textheight}
\end{question}
\begin{question}
\gkti{9}如图，正方体 $ABCD-A_1B_1C_1D_1$ 的棱线长为 $1$，线段 $B_1D_1$ 上有两个动点 $E$，$F$，且 $EF=\frac{\sqrt2}{2}$，则下列结论中错误的是（　）
\nopagebreak[4]
\begin{center}% 2009 年新课标卷文科第 9 题：正方体 ABCD-A1B1C1D1，E、F 在 B1D1 上。
% 来源 PDF：普通高考/2009/2009新课标文(琼,宁).pdf，
% SHA256 9727691ae735985e787f7ae3febf48dc9e6c6cdd1aa53b465f5d95648abdd627
% 原页：第 1 页（A3 横向）中部，印刷页 543。
% 重绘依据：正方体棱长 1；E、F 是上底面对角线 B1D1 上的两个动点，EF=√2/2；
% 虚线为被遮挡的棱及辅助线，实线为可见棱与 B1D1。
\begin{tikzpicture}[scale=0.92, semithick, >=Stealth,
  every node/.style={inner sep=1.5pt, font=\small}]
  % 下底面 ABCD，上底面 A1B1C1D1
  \coordinate (A)  at (0,0);
  \coordinate (B)  at (4,0);
  \coordinate (C)  at (4.8,2.0);
  \coordinate (D)  at (0.8,2.0);
  \coordinate (A1) at (0,2.8);
  \coordinate (B1) at (4,2.8);
  \coordinate (C1) at (4.8,4.8);
  \coordinate (D1) at (0.8,4.8);
  % 可见棱：下底面前两条、四条侧棱、上底面
  \draw (A) -- (B) -- (C);
  \draw (A) -- (A1);
  \draw (B) -- (B1);
  \draw (C) -- (C1);
  \draw (A1) -- (B1) -- (C1) -- (D1) -- cycle;
  % 被遮挡的棱
  \draw[dashed] (C) -- (D) -- (A);
  \draw[dashed] (D) -- (D1);
  % 上底面对角线 B1D1 与动点 E、F
  \draw (B1) -- (D1);
  \coordinate (E) at (2.72,3.60);
  \coordinate (F) at (1.76,4.20);
  \fill (E) circle (2.2pt) node[above left] {$E$};
  \fill (F) circle (2.2pt) node[above right] {$F$};
  % 辅助线（原图中的虚线）
  \draw[dashed] (A) -- (E);
  \draw[dashed] (B) -- (F);
  \draw[dashed] (B) -- (D);
  \draw[dashed] (A) -- (C);
  % 顶点标注
  \node[below left]  at (A)  {$A$};
  \node[below right] at (B)  {$B$};
  \node[right]       at (C)  {$C$};
  \node[left]        at (D)  {$D$};
  \node[left]        at (A1) {$A_1$};
  \node[right]       at (B1) {$B_1$};
  \node[above right] at (C1) {$C_1$};
  \node[above left]  at (D1) {$D_1$};
\end{tikzpicture}
\end{center}
\nopagebreak[4]
（A）$AC\perp BE$

（B）$EF\parallel$ 平面 $ABCD$

（C）三棱锥 $A-BEF$ 的体积为定值

（D）异面直线 $AE$，$BF$ 所成的角为定值

\Needspace{0.4\textheight}
\end{question}
\begin{question}
\gkti{10}如果执行如图的程序框图，输入 $x=-2$，$h=0.5$，那么输出的各个数的和等于（　）
\nopagebreak[4]
\begin{center}% 2009 年新课标卷文科第 10 题：分段函数求值并累加输出的程序框图。
% 来源 PDF：普通高考/2009/2009新课标文(琼,宁).pdf，
% SHA256 9727691ae735985e787f7ae3febf48dc9e6c6cdd1aa53b465f5d95648abdd627
% 原页：第 1 页（A3 横向）中部，印刷页 543。
% 重绘依据：开始 → 输入 x,h → 判断 x<0：是则 y=0；否则判断 x<1：是则 y=x，否则 y=1；
%   汇合后输出 y → 判断 x≥2：否则 x=x+h 返回 x<0 之前，是则结束。
\begin{tikzpicture}[scale=0.82, semithick, >=Stealth,
  every node/.style={inner sep=3pt, font=\small}]
  \node (start) at (0,0) [draw, rounded corners=4pt, minimum width=2.0cm, minimum height=0.6cm] {开始};
  \node (input) at (0,-1.25) [draw, trapezium, trapezium left angle=75, trapezium right angle=105,
    minimum width=3.6cm, minimum height=0.7cm] {输入 $x,\,h$};
  \node (c1) at (0,-2.85) [draw, diamond, aspect=2.2, minimum width=3.0cm, minimum height=1.0cm] {$x<0$};
  \node (y0) at (-3.1,-4.45) [draw, minimum width=1.7cm, minimum height=0.62cm] {$y=0$};
  \node (c2) at (2.5,-4.45) [draw, diamond, aspect=2.0, minimum width=2.5cm, minimum height=0.92cm] {$x<1$};
  \node (yx) at (1.0,-5.95) [draw, minimum width=1.7cm, minimum height=0.62cm] {$y=x$};
  \node (y1) at (4.6,-5.95) [draw, minimum width=1.7cm, minimum height=0.62cm] {$y=1$};
  \coordinate (merge) at (0.8,-7.0);
  \node (out) at (0.8,-7.55) [draw, trapezium, trapezium left angle=75, trapezium right angle=105,
    minimum width=2.6cm, minimum height=0.68cm] {输出 $y$};
  \node (c3) at (0.8,-9.15) [draw, diamond, aspect=2.2, minimum width=2.8cm, minimum height=1.0cm] {$x\geqslant2$};
  \node (inc) at (5.9,-9.15) [draw, minimum width=2.2cm, minimum height=0.62cm] {$x=x+h$};
  \node (end) at (0.8,-10.8) [draw, rounded corners=4pt, minimum width=2.0cm, minimum height=0.6cm] {结束};

  \draw[->] (start) -- (input);
  \draw[->] (input) -- (c1);
  \draw[->] (c1.west) -| node[pos=0.22, above] {是} (y0.north);
  \draw[->] (c1.east) -| node[pos=0.22, above] {否} (c2.north);
  \draw[->] (c2.south) -- node[right=2pt, pos=0.45] {是} (yx.north);
  \draw[->] (c2.east) -| node[pos=0.22, above] {否} (y1.north);
  \draw (y0.south) -- (y0.south |- merge) -- (merge);
  \draw (yx.south) -- (yx.south |- merge) -- (merge);
  \draw (y1.south) -- (y1.south |- merge) -- (merge);
  \draw[->] (merge) -- (out.north);
  \draw[->] (out) -- (c3.north);
  \draw[->] (c3.east) -- node[above, pos=0.45] {否} (inc.west);
  \draw[->] (inc.north) |- (5.9,-1.9) -- (0,-1.9) -- (c1.north);
  \draw[->] (c3.south) -- node[right=2pt, pos=0.45] {是} (end.north);
\end{tikzpicture}
\end{center}
\nopagebreak[4]
（A）$3$

（B）$3.5$

（C）$4$

（D）$4.5$

\Needspace{0.4\textheight}
\end{question}
\begin{question}
\gkti{11}一个棱锥的三视图如图，则该棱锥的全面积（单位：$\mathrm{cm}^2$）为（　）
\nopagebreak[4]
\begin{center}% 2009 年新课标卷文科第 11 题：棱锥的三视图。
% 来源 PDF：普通高考/2009/2009新课标文(琼,宁).pdf，
% SHA256 9727691ae735985e787f7ae3febf48dc9e6c6cdd1aa53b465f5d95648abdd627
% 原页：第 1 页（A3 横向）右栏上部，印刷页 543。
% 重绘依据：正视图与侧视图都是底边 6、高 4 的等腰三角形；俯视图是直角边长 6 的
% 等腰直角三角形，顶点（锥顶的投影）在斜边中点，图中标出其到直角边的距离 3。
% 尺寸单位 cm。
\begin{tikzpicture}[scale=0.9, semithick, >=Stealth,
  every node/.style={inner sep=1.5pt, font=\small}]
  % 正视图
  \begin{scope}
    \draw (0,0) -- (6,0) -- (3,4) -- cycle;
    \draw[dashed] (3,4) -- (3,0);
    \draw[<->] (3,4.45) -- (6,4.45);
    \node[above] at (4.5,4.45) {$3$};
    \draw[<->] (0,-0.55) -- (6,-0.55);
    \node[below] at (3,-0.55) {$6$};
    \draw[<->] (6.55,0) -- (6.55,4);
    \node[right] at (6.6,2) {$4$};
    \node[below] at (3,-1.2) {正（主）视图};
  \end{scope}
  % 侧视图
  \begin{scope}[shift={(10.2,0)}]
    \draw (0,0) -- (6,0) -- (3,4) -- cycle;
    \draw[dashed] (3,4) -- (3,0);
    \draw[<->] (3,4.45) -- (6,4.45);
    \node[above] at (4.5,4.45) {$3$};
    \draw[<->] (0,-0.55) -- (6,-0.55);
    \node[below] at (3,-0.55) {$6$};
    \draw[<->] (6.55,0) -- (6.55,4);
    \node[right] at (6.6,2) {$4$};
    \node[below] at (3,-1.2) {侧（左）视图};
  \end{scope}
  % 俯视图
  \begin{scope}[shift={(0,-7.2)}]
    \draw (0,0) -- (6,0) -- (0,6) -- cycle;
    \draw (0,0) -- (3,3);
    \draw (3,3) -- (6,3);
    \draw[<->] (6.55,0) -- (6.55,3);
    \node[right] at (6.6,1.5) {$3$};
    \draw[<->] (-0.55,0) -- (-0.55,6);
    \node[left] at (-0.6,3) {$6$};
    \draw[<->] (0,-0.55) -- (6,-0.55);
    \node[below] at (3,-0.55) {$6$};
    \node[below] at (3,-1.2) {俯视图};
  \end{scope}
  \node at (10.2,-3.2) {（单位：$\mathrm{cm}$）};
\end{tikzpicture}
\end{center}
\nopagebreak[4]
（A）$48+12\sqrt2$

（B）$48+24\sqrt2$

（C）$36+12\sqrt2$

（D）$36+24\sqrt2$

\end{question}
\begin{question}
\gkti{12}用 $\min\{a,b,c\}$ 表示 $a$，$b$，$c$ 三个数中的最小值。设 $f(x)=\min\{2^x,x+2,10-x\}$（$x\geqslant0$），则 $f(x)$ 的最大值为（　）

（A）$4$

（B）$5$

（C）$6$

（D）$7$

\end{question}
\gkpart{二、填空题}

\begin{question}
\gkti{13}曲线 $y=x\mathrm e^x+2x+1$ 在点 $(0,1)$ 处的切线方程为$\underline{\qquad\qquad}$。

\end{question}
\begin{question}
\gkti{14}已知抛物线 $C$ 的顶点坐标为原点，焦点在 $x$ 轴上，直线 $y=x$ 与抛物线 $C$ 交于 $A$，$B$ 两点，若 $P(2,2)$ 为 $AB$ 的中点，则抛物线 $C$ 的方程为$\underline{\qquad\qquad}$。

\end{question}
\begin{question}
\gkti{15}等比数列 $\{a_n\}$ 的公比 $q>0$，已知 $a_2=1$，$a_{n+2}+a_{n+1}=6a_n$，则 $\{a_n\}$ 的前 $4$ 项和 $S_4=$$\underline{\qquad\qquad}$。

\Needspace{0.4\textheight}
\end{question}
\begin{question}
\gkti{16}已知函数 $f(x)=2\sin(\omega x+\varphi)$ 的图象如图所示，则 $f\left(\frac{7\pi}{12}\right)=$$\underline{\qquad\qquad}$。
\nopagebreak[4]
\begin{center}% 2009 年新课标卷文科第 16 题：f(x)=2sin(ωx+φ) 的图象。
% 来源 PDF：普通高考/2009/2009新课标文(琼,宁).pdf，
% SHA256 9727691ae735985e787f7ae3febf48dc9e6c6cdd1aa53b465f5d95648abdd627
% 原页：第 1 页（A3 横向）右栏下部，印刷页 543。
% 重绘依据：图象自 x=π/4 的上升零点画到 x=5π/4 的零点，其间含两个最大值、一个最小值，
% 即 1.5 个周期，故 T=2π/3、ω=3、φ=−3π/4，f(x)=2sin(3x−3π/4)。
% 原页只画这一段（x<π/4 未画），本图照原页范围绘制，不补画左侧。
\begin{tikzpicture}[x=1.9cm, y=1.0cm, semithick, >=Stealth,
  every node/.style={inner sep=1.5pt, font=\small}]
  \draw[->] (-0.35,0) -- (4.55,0) node[right] {$x$};
  \draw[->] (0,-2.55) -- (0,2.55) node[above] {$y$};
  \draw (-0.09,1) -- (0.09,1) node[right=7pt] {$1$};
  \draw (-0.09,2) -- (0.09,2) node[right=7pt] {$2$};
  \draw (-0.09,-1) -- (0.09,-1) node[right=7pt] {$-1$};
  \draw (-0.09,-2) -- (0.09,-2) node[right=7pt] {$-2$};
  \node[below left] at (0,0) {$O$};
  % 曲线：f(x)=2sin(3x−3π/4)，x 从 π/4 到 5π/4
  \draw[domain=0.785398:3.926991, samples=320, smooth]
    plot (\x, {2*sin(deg(3*\x-2.356194))});
  % x 轴刻度标注（与原页一致，只标两个零点）
  \draw (0.785398,0) -- (0.785398,-0.12);
  \node[below] at (0.785398,-0.12) {$\frac{\pi}{4}$};
  \draw (3.926991,0) -- (3.926991,-0.12);
  \node[below] at (3.926991,-0.12) {$\frac{5\pi}{4}$};
\end{tikzpicture}
\end{center}
\nopagebreak[4]
\end{question}
\gkpart{三、解答题}

\Needspace{0.4\textheight}
\begin{question}
\gkti{17}如图，为了解某海域海底构造，在海平面内一条直线上的 $A$，$B$，$C$ 三点进行测量，已知 $AB=50\ \mathrm m$，$BC=120\ \mathrm m$，于 $A$ 处测得水深 $AD=80\ \mathrm m$，于 $B$ 处测得水深 $BE=200\ \mathrm m$，于 $C$ 处测得水深 $CF=110\ \mathrm m$，求 $\angle DEF$ 的余弦值。
\nopagebreak[4]
\begin{center}% 2009 年新课标卷文科第 17 题：海平面上的 A、B、C 与海底的 D、E、F。
% 来源 PDF：普通高考/2009/2009新课标文(琼,宁).pdf，
% SHA256 9727691ae735985e787f7ae3febf48dc9e6c6cdd1aa53b465f5d95648abdd627
% 原页：第 1 页（A3 横向）右栏下部，印刷页 543。
% 重绘依据：A、B、C 在同一条海平面直线上，AB=50 m、BC=120 m；AD、BE、CF 分别
% 垂直于海平面，AD=80 m、BE=200 m、CF=110 m；D、E、F 为海底测点，∠DEF 的顶点为 E。
% 横向按 0.05 cm/m、纵向按 0.035 cm/m 作图；海底轮廓为示意，不参与计算。
\begin{tikzpicture}[x=0.05cm, y=0.035cm, semithick, >=Stealth,
  every node/.style={inner sep=1.5pt, font=\small}]
  % 海平面（双波浪线）
  \draw[domain=-15:205, samples=120]
    plot (\x, {3*sin(deg(0.09*\x))});
  \draw[domain=-15:205, samples=120]
    plot (\x, {-3*sin(deg(0.09*\x))-7});
  % 海底轮廓（示意）
  \draw[line width=1.4pt, gray!70]
    (-18,-98) -- (0,-80) -- (12,-92) -- (18,-108) -- (30,-128)
    -- (38,-150) -- (42,-176) -- (50,-200) -- (62,-188) -- (72,-172)
    -- (86,-150) -- (102,-132) -- (120,-118) -- (140,-112) -- (170,-110)
    -- (190,-122) -- (200,-130);
  % 竖直测线
  \draw (0,0) -- (0,-80);
  \draw (50,0) -- (50,-200);
  \draw (170,0) -- (170,-110);
  % 海底三角形 DEF 的两边
  \draw (0,-80) -- (50,-200) -- (170,-110);
  % 海平面点
  \fill (0,0) circle (2.2pt) node[above] {$A$};
  \fill (50,0) circle (2.2pt) node[above] {$B$};
  \fill (170,0) circle (2.2pt) node[above] {$C$};
  % 海底点
  \fill (0,-80) circle (2.2pt) node[below left] {$D$};
  \fill (50,-200) circle (2.2pt) node[below] {$E$};
  \fill (170,-110) circle (2.2pt) node[right] {$F$};
  % 水平距离
  \draw[<->] (0,18) -- (50,18);
  \node[above] at (25,18) {$50$};
  \draw[<->] (50,18) -- (170,18);
  \node[above] at (110,18) {$120$};
  % 水深
  \node[left] at (-8,-42) {$80$};
  \node[left] at (42,-100) {$200$};
  \node[right] at (178,-55) {$110$};
\end{tikzpicture}
\end{center}
\nopagebreak[4]
\Needspace{0.4\textheight}
\end{question}
\begin{question}
\gkti{18}如图，在三棱锥 $P-ABC$ 中，$\triangle PAB$ 是等边三角形，$\angle PAC=\angle PBC=90^\circ$。

（1）求证：$AB\perp PC$；

（2）若 $PC=4$，且平面 $PAC\perp$ 平面 $PBC$，求三棱锥 $P-ABC$ 体积。
\nopagebreak[4]
\begin{center}% 2009 年新课标卷文科第 18 题：三棱锥 P-ABC。
% 来源 PDF：普通高考/2009/2009新课标文(琼,宁).pdf，
% SHA256 9727691ae735985e787f7ae3febf48dc9e6c6cdd1aa53b465f5d95648abdd627
% 原页：第 2 页（A3 横向）左上，印刷页 544。
% 重绘依据：△PAB 为等边三角形，∠PAC=∠PBC=90°，平面 PAC⊥平面 PBC。
% 按原图遮挡关系，PBC 为可见的前面（实线），A 在后方，PA、AB、AC 画虚线。
\begin{tikzpicture}[scale=0.95, semithick, >=Stealth,
  every node/.style={inner sep=1.5pt, font=\small}]
  \coordinate (P) at (1.0,5.0);
  \coordinate (B) at (1.8,0.4);
  \coordinate (C) at (8.2,1.2);
  \coordinate (A) at (4.7,2.3);
  % 可见棱
  \draw (P) -- (B) -- (C) -- (P);
  % 被遮挡的棱
  \draw[dashed] (P) -- (A);
  \draw[dashed] (A) -- (B);
  \draw[dashed] (A) -- (C);
  % 点与标注
  \fill (P) circle (2.2pt) node[above] {$P$};
  \fill (A) circle (2.0pt) node[below right] {$A$};
  \fill (B) circle (2.0pt) node[below left] {$B$};
  \fill (C) circle (2.0pt) node[right] {$C$};
\end{tikzpicture}
\end{center}
\nopagebreak[4]
\Needspace{0.4\textheight}
\end{question}
\begin{question}
\gkti{19}某工厂有工人 $1000$ 名，其中 $250$ 名工人参加过短期培训（称为 $A$ 类工人），另外 $750$ 名工人参加过长期培训（称为 $B$ 类工人），现用分层抽样方法（按 $A$ 类、$B$ 类分二层）从该工厂的工人中共抽查 $100$ 名工人，调查他们的生产能力（此处生产能力指一天加工的零件数）。

（1）求甲、乙两工人都被抽到的概率，其中甲为 $A$ 类工人，乙为 $B$ 类工人；

（2）从 $A$ 类工人中的抽查结果和从 $B$ 类工人中的抽插结果分别如下表 1 和表 2。

表 1：

\begin{center}
\begin{tabular}{llllll}
\toprule
生产能力分组 & $[100,110)$ & $[110,120)$ & $[120,130)$ & $[130,140)$ & $[140,150)$ \\
\midrule
人数 & $4$ & $8$ & $x$ & $5$ & $3$ \\
\bottomrule
\end{tabular}
\end{center}

表 2：

\begin{center}
\begin{tabular}{lllll}
\toprule
生产能力分组 & $[110,120)$ & $[120,130)$ & $[130,140)$ & $[140,150)$ \\
\midrule
人数 & $6$ & $y$ & $36$ & $18$ \\
\bottomrule
\end{tabular}
\end{center}

① 先确定 $x$，$y$，再在答题纸上完成下列频率分布直方图。就生产能力而言，$A$ 类工人中个体间的差异程度与 $B$ 类工人中个体间的差异程度哪个更小？（不用计算，可通过观察直方图直接回答结论）
\nopagebreak[4]
\begin{center}% 2009 年新课标卷文科第 19 题：A 类、B 类工人生产能力的频率分布直方图（空白网格）。
% 来源 PDF：普通高考/2009/2009新课标文(琼,宁).pdf，
% SHA256 9727691ae735985e787f7ae3febf48dc9e6c6cdd1aa53b465f5d95648abdd627
% 原页：第 2 页（A3 横向）左下，印刷页 544。
% 重绘依据：两组组距均为 10，横轴 100–150；A 类纵轴 0.004–0.036（步长 0.004），
% B 类纵轴 0.004–0.048（步长 0.004）。原页只给空白网格，矩形条由答题者补全。
\begin{tikzpicture}[scale=0.78, semithick, >=Stealth,
  every node/.style={inner sep=1.5pt, font=\small}]
  % A 类
  \begin{scope}
    \draw[->] (0,0) -- (5.9,0) node[right] {生产能力};
    \draw[->] (0,0) -- (0,4.9) node[above] {频率/组距};
    \foreach \i in {0,1,2,3,4,5} {
      \draw (\i,0) -- (\i,-0.14) node[below,font=\scriptsize] {\pgfmathparse{int(100+10*\i)}\pgfmathresult};
    }
    \foreach \k in {1,...,9} {
      \draw[gray!55] (0,0.5*\k) -- (5,0.5*\k);
      \draw (-0.12,0.5*\k) -- (0,0.5*\k)
        node[left] {\pgfmathparse{0.004*\k}\pgfmathprintnumber[fixed,precision=3]{\pgfmathresult}};
    }
    \foreach \i in {1,...,5} {
      \draw[gray!55] (\i,0) -- (\i,4.5);
    }
    \node[below] at (2.5,-2.0) {A 类工人生产能力的频率分布直方图};
  \end{scope}
  % B 类
  \begin{scope}[shift={(10.4,0)}]
    \draw[->] (0,0) -- (5.9,0) node[right] {生产能力};
    \draw[->] (0,0) -- (0,6.4) node[above] {频率/组距};
    \foreach \i in {0,1,2,3,4,5} {
      \draw (\i,0) -- (\i,-0.14) node[below,font=\scriptsize] {\pgfmathparse{int(100+10*\i)}\pgfmathresult};
    }
    \foreach \k in {1,...,12} {
      \draw[gray!55] (0,0.5*\k) -- (5,0.5*\k);
      \draw (-0.12,0.5*\k) -- (0,0.5*\k)
        node[left] {\pgfmathparse{0.004*\k}\pgfmathprintnumber[fixed,precision=3]{\pgfmathresult}};
    }
    \foreach \i in {1,...,5} {
      \draw[gray!55] (\i,0) -- (\i,6.0);
    }
    \node[below] at (2.5,-2.0) {B 类工人生产能力的频率分布直方图};
  \end{scope}
\end{tikzpicture}
\end{center}
\nopagebreak[4]
② 分别估计 $A$ 类工人和 $B$ 类工人生产能力的平均数，并估计该工厂工人的生产能力的平均数。（同一组中的数据用该组区间的中点值作代表）

\end{question}
\begin{question}
\gkti{20}已知椭圆 $C$ 的中心为直角坐标系 $xOy$ 的原点，焦点在 $x$ 轴上，它的一个顶点到两个焦点的距离分别是 $7$ 和 $1$。

（1）求椭圆 $C$ 的方程；

（2）若 $P$ 为椭圆 $C$ 上的动点，$M$ 为过 $P$ 且垂直于 $x$ 轴的直线上的点，$\frac{|OP|}{|OM|}=e$（$e$ 为椭圆 $C$ 的离心率），求点 $M$ 的轨迹方程，并说明轨迹是什么曲线。

\end{question}
\begin{question}
\gkti{21}已知函数 $f(x)=x^3-3ax^2-9a^2x+a^3$。

（1）设 $a=1$，求函数 $f(x)$ 的极值；

（2）若 $a>\frac14$，且当 $x\in[1,4a]$ 时，$|f'(x)|\leqslant12a$ 恒成立，试确定 $a$ 的取值范围。

\end{question}
\gkpart{四、选做题}

\Needspace{0.4\textheight}
\begin{question}
\gkti{22}如图，已知 $\triangle ABC$ 的两条角平分线 $AD$ 和 $CE$ 相交于 $H$，$\angle B=60^\circ$，$F$ 在 $AC$ 上，且 $AE=AF$。

（1）证明：证明：$B$，$D$，$H$，$E$ 四点共圆；

（2）证明：$CE$ 平分 $\angle DEF$。
\nopagebreak[4]
\begin{center}% 2009 年新课标卷文科第 22 题：△ABC 的两条角平分线 AD、CE 交于 H，F 在 AC 上且 AE=AF。
% 来源 PDF：普通高考/2009/2009新课标文(琼,宁).pdf，
% SHA256 9727691ae735985e787f7ae3febf48dc9e6c6cdd1aa53b465f5d95648abdd627
% 原页：第 2 页（A3 横向）右上，印刷页 544。
% 重绘依据：∠B=60°；AD 平分 ∠A，CE 平分 ∠C；E 在 AB 上、D 在 BC 上；
% F 在 AC 上且 AE=AF；H 为 AD 与 CE 的交点（内心）。图中按此几何关系取点。
\begin{tikzpicture}[scale=0.95, semithick, >=Stealth,
  every node/.style={inner sep=1.5pt, font=\small}]
  \coordinate (B) at (0,0);
  \coordinate (C) at (8,0);
  \coordinate (A) at (2.6,6.0);
  \coordinate (D) at (3.58,0);
  \coordinate (E) at (1.30,3.0);
  \coordinate (F) at (4.79,3.57);
  \coordinate (H) at (3.23,2.14);
  % 三角形
  \draw (A) -- (B) -- (C) -- cycle;
  % 角平分线
  \draw (A) -- (D);
  \draw (C) -- (E);
  % 辅助线 EF、ED
  \draw (E) -- (F);
  \draw (E) -- (D);
  % 点与标注
  \fill (A) circle (2.0pt) node[above] {$A$};
  \fill (B) circle (2.0pt) node[left] {$B$};
  \fill (C) circle (2.0pt) node[right] {$C$};
  \fill (D) circle (1.8pt) node[below] {$D$};
  \fill (E) circle (1.8pt) node[left] {$E$};
  \fill (F) circle (1.8pt) node[right] {$F$};
  \fill (H) circle (1.8pt) node[below left] {$H$};
\end{tikzpicture}
\end{center}
\nopagebreak[4]
\end{question}
\begin{question}
\gkti{23}已知曲线 $C_1:\begin{cases}x=-4+\cos t,\\ y=3+\sin t,\end{cases}$（$t$ 为参数），$C_2:\begin{cases}x=8\cos\theta,\\ y=3\sin\theta,\end{cases}$（$\theta$ 为参数）。

（1）化 $C_1$，$C_2$ 的方程为普通方程，并说明它们分别表示什么曲线；

\Needspace{5\baselineskip}
（2）若 $C_1$ 上的点 $P$ 对应的参数为 $t=\frac{\pi}{2}$，$Q$ 为 $C_2$ 上的动点，求 $PQ$ 中点 $M$ 到直线

\nopagebreak[4]
\[C_3:\begin{cases}x=3+2t,\\ y=-2+t,\end{cases}\quad(t\text{ \text{为参数}})\]

距离的最小值。

\Needspace{0.4\textheight}
\end{question}
\begin{question}
\gkti{24}如图，$O$ 为数轴的原点，$A$，$B$，$M$ 为数轴上三点，$C$ 为线段 $OM$ 上的动点，设 $x$ 表示 $C$ 与原点的距离，$y$ 表示 $C$ 到 $A$ 距离 $4$ 倍与 $C$ 到 $B$ 距离的 $6$ 倍的和。

（1）将 $y$ 表示成 $x$ 的函数；

（2）要使 $y$ 的值不超过 $70$，$x$ 应该在什么范围内取值？
\nopagebreak[4]
\begin{center}% 2009 年新课标卷文科第 24 题：数轴上的 O、A、B、M 与线段 OM 上的动点 C。
% 来源 PDF：普通高考/2009/2009新课标文(琼,宁).pdf，
% SHA256 9727691ae735985e787f7ae3febf48dc9e6c6cdd1aa53b465f5d95648abdd627
% 原页：第 2 页（A3 横向）右下，印刷页 544。
% 重绘依据：O 为原点，A、B、M 依次对应 10、20、30；C 在线段 OM 上，x=|OC|。
\begin{tikzpicture}[scale=0.9, semithick, >=Stealth,
  every node/.style={inner sep=1.5pt, font=\small}]
  \draw[->] (-1.2,0) -- (9.6,0);
  \foreach \x/\lab in {0/O,3/A,6/B,9/M} {
    \fill (\x,0) circle (2.0pt);
    \node[above] at (\x,0.02) {$\lab$};
  }
  \foreach \x/\lab in {3/10,6/20,9/30} {
    \draw (\x,0) -- (\x,-0.16);
    \node[below] at (\x,-0.18) {$\lab$};
  }
  \fill (5.0,0) circle (2.0pt) node[below] {$C$};
  \draw[<->] (0,0.55) -- (5.0,0.55);
  \node[above] at (2.5,0.55) {$x$};
\end{tikzpicture}
\end{center}
\nopagebreak[4]

\end{question}